% JAC1964-20261121. /root/arithmetic_sign; GPT-6 Astra Pro.
\subsection*{AR-20. Why finite-field Chow injectivity does not detect the arithmetic complement}

\begin{lemma}[AR-20a]\label{ar20a}
Let \(B/\mathbf F_q\) be an abelian variety, let \(h^j(B)\) be its
canonical Chow--K\"unneth summand, and let \(r\) be an integer.
Then
\[
     \operatorname{CH}^r(h^j(B))_{\mathbf Q}=0
                     \qquad\text{if }j\ne2r.                 \tag{AR20.1}
\]
This assertion uses the known finite-dimensionality of abelian
Chow motives and the Weil bounds; it does not require the Tate
conjecture.
\end{lemma}

\begin{proof}
Let \(F:B\to B\) be the absolute \(q\)-power Frobenius, which is a
morphism over \(\mathbf F_q\), and write \(F^*\) for its action on
the Chow motive. The canonical projectors commute with \(F^*\).
Let \(P_j(t)\in\mathbf Q[t]\) be the characteristic polynomial of
\(F^*\) on \(H^j(B_{\overline{\mathbf F}_q},\mathbf Q_\ell)\),
where \(\ell\nmid q\). This polynomial has rational coefficients:
it is the appropriate exterior-power polynomial of the Weil
polynomial on \(H^1\). Every complex root of \(P_j\) has absolute
value \(q^{j/2}\).

The endomorphism \(P_j(F^*)\) of \(h^j(B)\) is homologically
trivial, hence numerically trivial. Since \(h^j(B)\) is a
finite-dimensional Chow motive, Kimura nilpotence gives an integer
\(a>0\) such that
\[
                    P_j(F^*)^a=0
             \quad\text{in }\operatorname{End}(h^j(B)).
             \tag{AR20.2}
\]
On codimension-\(r\) cycles over \(\mathbf F_q\), flat pullback by
the absolute \(q\)-power Frobenius is multiplication by \(q^r\).
Indeed Frobenius fixes every underlying integral subvariety and
its pullback has multiplicity \(q^r\) in codimension \(r\).
Thus for \(z\in\operatorname{CH}^r(h^j(B))_{\mathbf Q}\),
(AR20.2) gives
\[
                        P_j(q^r)^a z=0.
\]
When \(j\ne2r\), the Weil absolute values show that
\(P_j(q^r)\ne0\). It is an invertible rational scalar, so \(z=0\).
\end{proof}

\begin{corollary}[AR-20b]\label{ar20b}
Let \(J/K_0\) be the Fermat Jacobian, with
\(K_0=\mathbf Q(i)\), and let \(N\subset h^3(J)\) be the
canonical complement. At every finite place of good reduction
where these projectors are spread out, the specialization map
\[
 \operatorname{sp}_v:
 \operatorname{CH}^2(N_{K_0})_{\mathbf Q}
       \longrightarrow
 \operatorname{CH}^2(N_{k(v)})_{\mathbf Q}
                                                               \tag{AR20.3}
\]
is the zero map, since its target is zero. The same holds after
any finite extension of the number field.
In particular the nonzero, positive-height classes
\(\eta,T\eta\) of AR-13 and JAC-P11 specialize to zero rationally
at all such places.
\end{corollary}

\begin{proof}
Specialization for a smooth proper model respects correspondence
actions and the canonical abelian Chow--K\"unneth projectors.
It therefore takes the displayed group into
\(\operatorname{CH}^2(h^3(J_{k(v)}))_{\mathbf Q}\), which is zero
by AR-20a with \(j=3\) and \(r=2\). If a chosen correspondence
requires a finite residue extension, apply the result there and
use restriction--corestriction on rational Chow groups.
The global classes \(\eta,T\eta\) are nonzero by their positive
height matrix. Thus they are explicit nonzero elements in the
kernel of the product of all these specialization maps.
\end{proof}

\paragraph{Exact limitation of the specialization route.}
JAC-P15's finite-field rational--numerical comparison cannot
be promoted to injectivity of the arithmetic Abel--Jacobi map
by testing these finite fibres: all of the arithmetic complement
already disappears in their rational Chow groups. No failure of
JAC-P15 is involved. The missing information concerns the
number-field extension class and its global regulator.

This limitation also appears in the general cycle-class
comparison. Kahn, \emph{On the injectivity and non-injectivity of
the \(l\)-adic cycle class maps}, arXiv:2307.13281v2
(10 April 2024), Section~4, printed pp.~10--12, explains that
passing from Chow groups to completed \'{e}tale motivic Chow
groups may acquire kernel even in an attempted argument which
assumes finite generation of ordinary Chow groups. His
Theorem~5.1, printed p.~13, derives the positive-characteristic
statement from the full Tate--Beilinson conjecture, while
Section~5.7, printed pp.~17--18, does not give a
characteristic-zero analogue. The decomposition-of-the-diagonal
criterion in Theorem~6.1, printed p.~18, assumes
\(\operatorname{CH}_0(X_{k(X)})_{\mathbf Q}=\mathbf Q\);
it is not available for \(J\). In particular a factorization of
the whole summand \(N(1)\) through degree-one motives would
contradict its Hodge types \((2,-1),(-1,2)\).

\paragraph{Established inputs and check scope.}
The nilpotence mechanism is the same one explicitly proved in
Kahn 2003, Proposition~1.5 and Lemma~1.6, printed p.~981, and
used in JAC-P13 and JAC-P15. AR-20 applies that mechanism to
the actual Frobenius polynomial before any Tate-conjecture
input. The Chow--K\"unneth decomposition and its functoriality
are the established abelian-motive inputs already used in
AR-1--AR-2. No claim about a nonzero zero-height class or a
negative-height class is made here.
